\originalchapter{8.05 2013测试：尺度变换与测量宽度}
原件为Fall 2013 Test, 2013年10月23日12:30--14:00, 五题100分, 闭卷禁计算器. 以下为\textbf{AI独立解答}.
\originalsection{第1题：十个判断（20分）}
\begin{exercise}
判断: (1)$\Delta Q=0$必为$Q$本征态; (2)时间无显式依赖算符的宽度恒定; (3)$N\geq2$时可有有限$X,P$满足$[X,P]=\ii\hbar I$; (4)两个对易Hermitian算符中$Q_1$的每个本征矢都是$Q_2$本征矢. 后六项为有限维复空间: (5)恒等算符矩阵依赖基; (6)谱与基无关; (7)任何基中伴随矩阵都为共轭转置; (8)向真低维子空间的正交投影在全空间中既非单射又非满射; (9)正交归一组线性独立; (10)幺正且Hermitian者平方为恒等.
\end{exercise}
\begin{solution}
依次\textbf{T,F,F,F,F,T,F,T,T,T}. (1)$(Q-\langle Q\rangle)\ket\psi$范数为0; (2)态仍演化, 除非守恒或特定态; (3)对易子迹0而右边非零; (4)简并子空间内任意矢未必对角化另一算符, 例如$Q_1=I,Q_2=\sigma_z$; (5)每个基中$I$矩阵仍$I$; (6)相似变换不改谱; (7)仅正交归一基中直接成立, 一般基有Gram矩阵; (8)核非零、像非全空间; (9)内积逐项取系数; (10)$U^\dagger=U=U^{-1}$.
\end{solution}
\originalsection{第2题：幺正尺度变换（15分）}
\begin{exercise}
$U_\lambda=\exp[-\ii\lambda(xp+px)/(2\hbar)]$. (a)幺正要求$\lambda$何条件. (b)求$U_\lambda^\dagger xU_\lambda$. (c)求动量变换, 判断是否与自由$H$对易.
\end{exercise}
\begin{solution}
正规自伴伸缩生成元$D=(xp+px)/2$要求实$\lambda$使指数反Hermitian. $[D,x]=-\ii\hbar x$,$[D,p]=\ii\hbar p$, 令$X(\lambda)=U_\lambda^\dagger xU_\lambda$, 对参数求导得$X\prime=(\ii/\hbar)U^\dagger[D,x]U=X$, 初值$x$给$X=\e^\lambda x$. 同理$P\prime=-P$, 得$U^\dagger pU=\e^{-\lambda}p$. 位置波包拉宽$\e^\lambda$时动量宽度缩小$\e^{-\lambda}$, 总概率保持. 对易$[\e^\lambda x,\e^{-\lambda}p]=\ii\hbar$保持, 但自由$H$变为$\e^{-2\lambda}H$, 一般不相等, 所以除$\lambda=0$外不是自由Hamiltonian的对称性. 幺正等谱的陈述与缩放动量能量并不矛盾, 因自由连续谱$[0,\infty)$在正缩放下集合相同.
\end{solution}
\originalsection{第3题：由长度恢复内积（15分）}
\begin{exercise}
(a)化简$\tfrac14[\norm{u+v}^2-\norm{u-v}^2+\ii\norm{u-\ii v}^2-\ii\norm{u+\ii v}^2]$, 说明保范数的线性$S$保内积. (b)从该式及$\norm{\alpha u}=|\alpha|\norm u$证明共轭对称.
\end{exercise}
\begin{solution}
展开得$\inner uv$（本册内积对第二项线性）. 线性保证$S(u\pm v)=Su\pm Sv$以及带$\ii$组合相同, 四个范数保持就使整个极化式保持. 交换$u,v$, 实部分$\norm{v-u}=\norm{u-v}$不变. 用$\norm{v-\ii u}=\norm{-\ii(u+\ii v)}=\norm{u+\ii v}$, $\norm{v+\ii u}=\norm{\ii(u-\ii v)}$使虚部交换符号, 因而$\inner vu=\inner uv^*$. 这里“内积由范数产生”包含极化式确实具有线性等内积性质的前提; 不是任意Banach范数都给内积.
\end{solution}
\originalsection{第4题：双阱势的变分估计（25分）}
\begin{exercise}
$H=p^2/(2m)+\hbar\omega-m\omega^2x^2/2+m^2\omega^3x^4/(16\hbar)$,$d_0^2=\hbar/(m\omega)$. (a)画势与极值. (b)对试探态$\psi=x\e^{-x^2/(2\lambda^2)}$求能量中四次项系数$A$. 已给$E=\hbar\omega-3m\omega^2\lambda^2/4+A\lambda^4+3\hbar^2/(4m\lambda^2)$. (c)令$\lambda=d_0\eta$, 忽略最后动能项优化再代回给上界. (d)这对哪一另一能级给更强界.
\end{exercise}
\begin{solution}
完成平方$V=m^2\omega^3(x^2-4d_0^2)^2/(16\hbar)$. 极小$x=\pm2d_0,V=0$, 局部极大$x=0,V=\hbar\omega$. 本题原卷积分常数约定$\int u^{2n}\e^{-u^2}\dd u=c_n\sqrt\pi$, $c_1=1/2,c_3=15/8$, 与上一年测试的$\e^{-2y^2}$常数不同. 试探态密度含$x^2$, 所以$\langle x^4\rangle=\lambda^4(c_3/c_1)=15\lambda^4/4$, $A=15m^2\omega^3/(64\hbar)$. 因而
\[
E(\eta)=\hbar\omega\left[1-\frac34\eta^2+\frac{15}{64}\eta^4+\frac3{4\eta^2}\right].
\]
忽略最后项只用于\textbf{选取试探参数}, 极值$\eta^2=8/5$. 将其代回完整能量而非删项能量得$E_*=139\hbar\omega/160=0.86875\hbar\omega$. 任一归一试探态的完整能量是基态上界; 删掉正动能项后的最小值并不是上界. 本试探态奇, 正交于偶基态, 展开只含奇能级, 故更直接$E_1\leq E_*$, 为第一激发态上界, 也当然高于真正基态.
\begin{center}\begin{tikzpicture}\begin{axis}[width=.72\textwidth,height=4cm,xlabel={$x/d_0$},ylabel={$V/(\hbar\omega)$},xmin=-3,xmax=3,ymin=0,ymax=2]
\addplot[main,domain=-3:3,samples=100]{(x*x-4)^2/16};\end{axis}\end{tikzpicture}\end{center}
\end{solution}
\originalsection{第5题：自旋宽度的上下限（25分）}
\begin{exercise}
自旋$1/2$: (a)推$\Delta S_x\Delta S_y$不等式. (b)最小乘积与取到的态. (c)任一分量最大宽度. (d)对一般$\ket{\boldsymbol n;+}$求三宽度和使$x,y$乘积最大的态.
\end{exercise}
\begin{solution}
$[S_x,S_y]=\ii\hbar S_z$给$\Delta S_x\Delta S_y\geq(\hbar/2)|\langle S_z\rangle|$. $S_i^2=\hbar^2I/4$, $\langle S_i\rangle=(\hbar/2)n_i$, 所以$\Delta S_i=(\hbar/2)\sqrt{1-n_i^2}$. 最小乘积0在$x$本征态$\ket{x;\pm}$或$y$本征态$\ket{y;\pm}$取到. 每个分量宽度上限$\hbar/2$, 当$n_i=0$取到. 一般角表示$n_x=\sin\theta\cos\phi,n_y=\sin\theta\sin\phi,n_z=\cos\theta$, 代入即三个公式. 最大乘积$\hbar^2/4$要求$n_x=n_y=0$, 即$\ket{z;\pm}$; 在这些态Robertson界也饱和. 测量宽度由制备决定, 不表示实验仪器的读数误差.
\end{solution}
